% Standalone problem document, generated from the adjacent README.md.
% Compile with XeLaTeX. Mathematical content is maintained in Markdown.
\documentclass[11pt,a4paper]{article}
\usepackage[fontsize=10.5pt]{scrextend}
\usepackage[margin=22mm,headheight=15pt,headsep=9mm,footskip=11mm]{geometry}
\usepackage{fontspec}
\setmainfont{texgyrepagella}[Extension=.otf,UprightFont=*-regular,BoldFont=*-bold,ItalicFont=*-italic,BoldItalicFont=*-bolditalic]
\setsansfont{texgyreheros}[Extension=.otf,UprightFont=*-regular,BoldFont=*-bold,ItalicFont=*-italic,BoldItalicFont=*-bolditalic]
\setmonofont{texgyrecursor}[Extension=.otf,UprightFont=*-regular,BoldFont=*-bold,ItalicFont=*-italic,BoldItalicFont=*-bolditalic,Scale=MatchLowercase]
\usepackage{amsmath,amssymb}
\usepackage{unicode-math}
\setmathfont{texgyrepagella-math.otf}
\usepackage{xcolor,microtype,enumitem,needspace,fancyhdr}
\usepackage{bookmark}
\definecolor{ink}{HTML}{17344B}
\definecolor{accent}{HTML}{197D80}
\definecolor{muted}{HTML}{596773}
\hypersetup{colorlinks=true,linkcolor=accent,urlcolor=accent,pdftitle={MD-06 - Global
synchronization of a random cubic
graph},pdfauthor={Open Problems in Numerical Linear Algebra}}
\urlstyle{same}
\setlength{\parindent}{0pt}
\setlength{\parskip}{5pt plus 1pt minus 1pt}
\linespread{1.04}
\setlength{\emergencystretch}{3em}
\widowpenalty=10000
\clubpenalty=10000
\displaywidowpenalty=10000
\allowdisplaybreaks[1]
\setlist{nosep,leftmargin=1.5em}
\providecommand{\tightlist}{\setlength{\itemsep}{0pt}\setlength{\parskip}{0pt}}
\providecommand{\pandocbounded}[1]{#1}
\pagestyle{fancy}
\fancyhf{}
\fancyhead[L]{\sffamily\footnotesize\color{muted}OPEN PROBLEMS IN NUMERICAL LINEAR ALGEBRA}
\fancyhead[R]{\sffamily\bfseries\footnotesize\color{accent}MD-06}
\fancyfoot[L]{\parbox[b]{0.9\textwidth}{\sffamily\fontsize{8}{10}\selectfont\color{muted}Editorial ratings; a bounded search cannot certify that a problem remains unsolved.\\Literature check: 2026-09-10}}
\fancyfoot[R]{\sffamily\footnotesize\color{muted}\thepage}
\renewcommand{\headrulewidth}{0.3pt}
\renewcommand{\headrule}{\hbox to\headwidth{\color{accent}\leaders\hrule height \headrulewidth\hfill}}
\makeatletter
\renewcommand{\section}{\Needspace{5\baselineskip}\@startsection{section}{1}{0pt}{12pt plus 2pt minus 2pt}{4pt}{\sffamily\bfseries\large\color{ink}}}
\renewcommand{\subsection}{\Needspace{4\baselineskip}\@startsection{subsection}{2}{0pt}{9pt plus 1pt minus 1pt}{3pt}{\sffamily\bfseries\normalsize\color{ink}}}
\makeatother
\setcounter{secnumdepth}{0}
\begin{document}
{\sffamily\small\color{muted}Matrix discrepancy and optimization\par}
\vspace{3pt}
{\raggedright\hyphenpenalty=10000\exhyphenpenalty=10000\sffamily\bfseries\fontsize{21}{24}\selectfont\color{ink}Global
synchronization of a random cubic graph\par}
\vspace{7pt}
{\sffamily\small\textbf{Difficulty:} challenging\quad\textbf{Importance:} interesting
to the community\par}
{\sffamily\small\bfseries\color{accent}Status: Open\par}
\vspace{4pt}
{\color{accent}\hrule height 0.8pt}
\vspace{6pt}
\textbf{Provenance:} source-stated conjecture, with the phase domain
written explicitly as a torus.

\textbf{Rating rationale:} Controlling every local minimum on sparse
random cubic graphs requires new landscape analysis beyond existing
dense and high-degree results; it informs synchronization and nonconvex
optimization.

For each even integer \(n\geq4\), let \(G_n\) be uniformly distributed
over the labelled simple \(3\)-regular graphs with vertex set
\(\{1,\ldots,n\}\). For a graph \(G\) on these vertices, define \[
E_G(\theta)=
\sum_{\{i,j\}\in E(G)}
\bigl(1-\cos(\theta_i-\theta_j)\bigr),
\qquad
\theta\in(\mathbb R/2\pi\mathbb Z)^n.
\] A phase vector is synchronized if \(\theta_i=\theta_j\) modulo
\(2\pi\) for every pair \(i,j\). A local minimum is understood with
respect to the usual topology on the phase torus; it need not be strict.

Is it true that \[
\lim_{\substack{n\to\infty\\ n\ {\rm even}}}
\mathbb P\!\left(
\text{every local minimum of }E_{G_n}
\text{ is synchronized}
\right)=1?
\]

The energy is a sparse graph optimization problem whose derivatives
involve weighted graph Laplacians. The conjecture asks whether a typical
graph of the smallest plausible fixed regular degree has no
nonsynchronized local minima.

\subsection{References}\label{references}

\begin{enumerate}
\def\labelenumi{\arabic{enumi}.}
\tightlist
\item
  A. S. Bandeira, A. Kireeva, A. Maillard, and A. Rödder,
  \emph{Randomstrasse101: Open Problems of 2024}, arXiv:2504.20539
  (2025), Definition 2.1 and Conjecture 3.
  \href{https://arxiv.org/html/2504.20539v1}{Paper}.
\item
  P. Abdalla, A. S. Bandeira, M. Kassabov, V. Souza, S. H. Strogatz, and
  A. Townsend, \emph{Expander graphs are globally synchronizing},
  Advances in Mathematics 488 (2026), article 110773;
  arXiv:2210.12788v4, introduction and Section 7. The random regular
  result requires substantially larger fixed degree.
  \href{https://arxiv.org/abs/2210.12788}{Paper}.
\item
  V. Jain, C. Mizgerd, and M. Sawhney, \emph{The random graph process is
  globally synchronizing}, arXiv:2501.12205 (2025), Theorem 1.2: the
  connectivity hitting-time result for a different random graph model.
  \href{https://arxiv.org/abs/2501.12205}{Paper}.
\item
  S. Lepsveridze and S. Zhang, \emph{Graphs with connectivity
  \(3/4-\varepsilon\) are globally synchronizing}, (2026), dense
  minimum-degree result. \href{https://arxiv.org/abs/2608.20010}{Paper}.
\end{enumerate}

\subsection{Status check --- 2026-09-10}\label{status-check-2026-09-10}

checked the current records 2504.20539v1 (2025-04-29), 2210.12788v4
(2026-01-19), 2501.12205v1 (2025-01-21), and 2608.20010v1 (2026-08-20),
and searched ``random cubic globally synchronizing 2026'' and ``random
regular graphs degree 3 conjecture''. The large-degree,
random-graph-process, and dense minimum-degree results do not establish
the stated cubic-graph assertion. No proof, counterexample, or relevant
withdrawal was found. The source's angular variables give the torus used
here; its isolated sphere notation is not imposed on the phases.

\textbf{Audit update (2026-09-10):} Rechecked the 2025 cubic-graph
conjecture and August 2026 dense-degree theorem, and searched for later
cubic synchronization results. The dense threshold result concerns a
different graph regime and does not resolve this target. This is a
bounded literature check, not a proof that no solution exists.
\end{document}
